% Source-aligned English for printed p.159.
\SourcePage{159}

\LayerTarget{tohoku:def:flasque-and-soft-sheaves}
A sheaf of sets $F$ on a space $X$ is called \textit{flabby} (respectively
\textit{soft}) if, for every open (respectively closed) subset $A$ of $X$,
every section of $F$ over $A$ is the restriction of a section over $X$. Thus,
given a set $E_x$ for every $x\in X$, the sheaf $E$ whose set of sections on
an open $U$ is
\[
E(U)=\prod_{x\in U}E_x
\]
with the evident restriction maps---compare the example before Proposition
3.1.2---is both flabby and soft. Hence every sheaf of sets embeds in a flabby
sheaf of sets, and every sheaf of abelian groups, or of $\mathbf O$-modules
when $\mathbf O$ is a sheaf of rings on $X$, embeds in a flabby sheaf of the
same kind. If a closed subset $A$ of $X$ has a paracompact neighborhood, then
every section over $A$ of a sheaf $F$ on $X$ extends to a section over some
neighborhood of $A$. It follows at once that \textit{every flabby sheaf on a
paracompact space $X$ is soft}.

\LayerTarget{tohoku:def:phi-soft-and-paracompactifying}
Now let $\Phi$ be a family of closed subsets of $X$ satisfying the general
conditions of 3.2. A sheaf $F$ of abelian groups on $X$ is
$\Phi$-\textit{soft} if, for every $A\in\Phi$ and every section of $F$ over
$A$, there is a section $f\in\Gamma_\Phi(F)$ inducing the given section. We
call $\Phi$ a \textit{paracompactifying family} if, in addition to the
conditions already imposed, it satisfies the following conditions, first
introduced in [4]: \textit{every $A\in\Phi$ is paracompact and has a
neighborhood $B\in\Phi$}. As above, one easily sees that \textit{a flabby
sheaf of abelian groups is $\Phi$-soft for every paracompactifying family
$\Phi$}. In particular, every sheaf of abelian groups embeds in a
$\Phi$-soft sheaf, since it even embeds in a flabby one. The value of these
definitions lies in the following

\LayerTarget{tohoku:prop:3.3.2}
\textsc{Proposition 3.3.2.} --- \textit{Let $X$ be a topological space
equipped with a family $\Phi$. Consider the functor $\Gamma_\Phi$ from the
category $C^X$ of sheaves of abelian groups on $X$ to the category of abelian
groups. The conditions of Lemma 3.3.1 hold in each of the following cases:}

\textit{1) $\mathbf M$ is the class of flabby sheaves on $X$.}

\textit{2) $\Phi$ is paracompactifying and $\mathbf M$ is the class of
$\Phi$-soft sheaves on $X$.}

\LayerTarget{tohoku:cor:3.3.2}
\textsc{Corollary.} --- \textit{$H^p_\Phi(X,F)=0$ for $p>0$ if $F$ is
flabby, or if $F$ is $\Phi$-soft and $\Phi$ is paracompactifying.}

Condition (i) of the lemma has already been verified, and condition (ii) is
immediate. Only condition (iii) requires a proof; we refer to Godement's book
for it, or offer it to the reader as an exercise.

\LayerTarget{tohoku:remark:fine-sheaves-are-soft}
\textsc{Remark.} --- If $\Phi$ is paracompactifying, one easily checks that
fine sheaves [4, XV] are $\Phi$-soft and therefore satisfy
$H^p_\Phi(X,F)=0$. Thus the cohomology theory given in [4] for
paracompactifying families $\Phi$ is a special case of the theory developed
here. We also refer to Godement's book for a particularly elegant definition

\TohokuContinuation{160}
